% !TeX root = ../../tesis.tex

% =============================================================================
%  4.2.5 — MOTOR DE IMPEDANCIA
% =============================================================================
\subsection{Motor de Impedancia: La Ecuación VFT}
\label{sec:vft-impedancia}

Una vez construido el grafo topológico, el módulo \texttt{VFTImpedanceModel}
recorre cada arista y le asigna un \textbf{peso temporal en minutos}: el tiempo
estimado de recorrido que los algoritmos de caminos mínimos utilizarán para
determinar la ruta más eficiente entre dos estaciones.

\subsubsection{Distancia Geodésica: Fórmula de Haversine}
\label{subsec:vft-haversine}

Todos los cálculos de distancia del modelo utilizan la
\textbf{fórmula de Haversine}, que obtiene la longitud del arco más corto sobre
la superficie de una esfera entre dos puntos de latitud y longitud conocidos. A
diferencia de la distancia euclidiana plana, Haversine opera sobre coordenadas
geográficas reales y produce metros con un error menor al 0.3~\% para distancias
urbanas dentro de la \zmvm{}.

La fórmula se calcula en dos pasos:

\begin{equation}
    a = \sin^{2}\!\left(\frac{\Delta\varphi}{2}\right)
        + \cos\varphi_1 \cdot \cos\varphi_2
          \cdot \sin^{2}\!\left(\frac{\Delta\lambda}{2}\right)
    \label{eq:haversine-a}
\end{equation}

\begin{equation}
    d = 2R \cdot \arctan2\!\left(\sqrt{a},\;\sqrt{1-a}\right)
    \label{eq:haversine-d}
\end{equation}

\begin{itemize}
    \item $\varphi_1, \varphi_2$: latitudes de los dos puntos en radianes.
    \item $\Delta\varphi = \varphi_2 - \varphi_1$: diferencia de latitudes.
    \item $\Delta\lambda = \lambda_2 - \lambda_1$: diferencia de longitudes.
    \item $R = 6{,}371{,}000$~m: radio medio de la Tierra.
    \item $d$: distancia geodésica resultante en metros.
\end{itemize}

La implementación como método estático de \texttt{VFTImpedanceModel} es
reutilizada por prácticamente todos los módulos del sistema. El código fuente
completo se presenta en el Anexo~\ref{anx:vft-codigo}.

\subsubsection{Coeficiente de Fricción Vial}
\label{subsec:vft-cf}

El \textbf{Coeficiente de Fricción Vial} ($CF$) traduce el tipo de derecho de
vía de cada segmento en un penalizador multiplicativo sobre el tiempo de
traslado. Se obtiene mediante:

\begin{equation}
    CF = 1 + \alpha \cdot \beta_{\text{ZMVM}}
    \label{eq:vft-cf}
\end{equation}

\begin{itemize}
    \item $\alpha \in \{0.0, 0.2, 0.5, 1.0\}$: factor de exposición al tráfico
    según el tipo de derecho de vía, con referencia normativa en
    \citep{sedatu2019manual}.
    \item $\beta_{\text{ZMVM}} = 0.759$: índice de saturación vial de la
    \zmvm{} \citep{tomtom2024}.
\end{itemize}

La clase \texttt{FrictionCalculator} implementa este cálculo de forma aislada,
aplicando vía mixta ($CF = 1.759$) como caso conservador por defecto. La
implementación Python completa se presenta en el Anexo~\ref{anx:vft-codigo}.

\subsubsection{Ecuación VFT: Peso Temporal de Cada Arista}
\label{subsec:vft-ecuacion}

Con la distancia geodésica y el coeficiente de fricción disponibles, el método
\path{apply_impedance()} asigna el peso temporal a cada arista del grafo:

\begin{equation}
    w_{ij} = \frac{d_{ij}}{v_{ij}} \cdot CF_{ij}
    \label{eq:vft-peso}
\end{equation}

\begin{itemize}
    \item $w_{ij}$: peso de la arista en minutos (tiempo de viaje estimado entre
    estaciones $i$ y $j$).
    \item $d_{ij}$: distancia real del trazo entre $i$ y $j$ en metros,
    acumulada vértice a vértice sobre el \texttt{MultiLineString} durante la
    construcción del grafo.
    \item $v_{ij}$: velocidad comercial del sistema en metros por minuto:
    $v_{ij} = v_{\text{km/h}} \times (1000/60)$.
    \item $CF_{ij}$: Coeficiente de Fricción Vial de la arista
    (Ecuación~\ref{eq:vft-cf}).
\end{itemize}

La implementación Python completa se presenta en el Anexo~\ref{anx:vft-codigo}.

Las aristas de transbordo peatonal quedan excluidas de este proceso porque ya
recibieron su peso durante la construcción del grafo: su costo suma el tiempo de
caminata entre estaciones más la espera promedio hasta el siguiente servicio
(véase la Sección~\ref{subsec:fase2-tb} del Capítulo~
\ref{ch:analisis-red-actual}). El peso de cada arista —en minutos— es la entrada de
todos los algoritmos de caminos mínimos ejecutados en los indicadores de Fase~3.
